\documentclass[11pt]{article}
\usepackage{mathunal-base}
\mathunalfuente{serif}

\renewcommand{\examtitulo}{Cálculo en Varias Variables (1000006) --- Segundo Examen Parcial}
\renewcommand{\examfecha}{Semestre 01-2013, 11 de Mayo de 2013}

\begin{document}

\examheader

\vspace{4pt}
\noindent Nombre: \dotfill\ Cédula: \dotfill

\vspace{4pt}
\noindent Profesor: \dotfill\ Grupo: \dotfill

\vspace{4pt}
\noindent Calificación: 1 \dotfill; 2 \dotfill; 3 \dotfill; 4 \dotfill. Total: \dotfill

\vspace{6pt}
\noindent\textbf{Instrucciones.} No está permitido el uso de calculadora ni de cualquier otro tipo de aparato electrónico, incluyendo teléfonos móviles, los cuales deben permanecer apagados durante el tiempo de duración del examen. Tampoco está permitido hacer preguntas a las personas encargadas de la vigilancia.

\vspace{6pt}
\noindent Duración: 1 hora y 50 minutos.

\vspace{4pt}
\noindent El examen consta de 4 puntos distribuidos en 4 páginas.

\begin{enumerate}
\item{} (25\%) Halle el valor máximo y el valor mínimo absolutos de la función
\[
f(x,y) = x+y^2
\]
en el conjunto cerrado y acotado
\[
D = \left\{(x,y)\ |\ 2x^2+y^2\leq 1\right\}.
\]
\noindent \textbf{Sugerencia}: Para hallar los valores extremos de $f$ en la frontera de $D$ use multiplicadores de Lagrange.

\newpage
\item{} (25\%) Halle la masa de una lámina delgada que ocupa la región $D$ en el primer cuadrante limitada por la parábola
\[
y = x^2
\]
y las rectas
\[
x=0 \quad \text{y} \quad y=4,
\]
si se sabe que la densidad en el punto $(x,y)$ es $\rho(x,y)=xe^{y^2}$.

\newpage
\item{} (25\%) Dibuje el sólido acotado por el paraboloide
\[
z = x^2+y^2
\]
y por la hoja superior $(z>0)$ del hiperboloide de dos hojas
\[
z^2-(x^2+y^2)=2
\]
y calcule su volumen.

\newpage
\item{} (25\%) Sea $\Omega$ el sólido definido por las condiciones
\[
\sqrt{3(x^2+y^2)}\leq z \quad \text{y} \quad x^2+y^2+z^2\leq 1.
\]
\noindent Sabiendo que la densidad en el punto $(x,y,z)$ del sólido está dada por
\[
\sigma(x,y,z) = 1+x^2+y^2+z^2,
\]
halle la masa del sólido y la coordenada $\bar{z}$ de su centro de masa.
\end{enumerate}

\examfooter
\end{document}
